\documentclass[10pt,twocolumn,letterpaper]{article}

% ============================================================
% Packages
% ============================================================
\usepackage{ctex}
\usepackage[T1]{fontenc}
\usepackage{amsmath,amssymb,amsfonts}
\usepackage{graphicx}
\usepackage{booktabs}
\usepackage[numbers,sort&compress]{natbib}
\usepackage{hyperref}
\usepackage{algorithm}
\usepackage{algorithmic}
\usepackage{multirow}
\usepackage{xcolor}
\usepackage{caption}
\usepackage{subcaption}
\usepackage{microtype}
\usepackage{enumitem}
\usepackage{bm}
\usepackage{mathtools}

% ============================================================
% Custom commands
% ============================================================
\newcommand{\method}{VITA-FASA}
\newcommand{\methodfull}{VITA with Flow-Aware Spatial Attention}
\newcommand{\fasa}{FASA}
\newcommand{\fasafull}{Flow-Aware Spatial Attention}
\newcommand{\RR}{\mathbb{R}}
\newcommand{\EE}{\mathbb{E}}
\newcommand{\LL}{\mathcal{L}}
\newcommand{\FF}{\mathcal{F}}
\newcommand{\MM}{\mathcal{M}}
\DeclareMathOperator*{\argmin}{arg\,min}

\title{看见何处：面向流匹配机器人策略的\\自监督空间注意力}

\author{
  匿名作者
}

\date{}

\begin{document}
\maketitle

% === 摘要 ===
\begin{abstract}
流匹配已成为视觉运动机器人策略的强大生成式框架，具有训练稳定和高质量多模态动作生成等优势。然而，现有的流匹配策略在视觉编码阶段采用全局平均池化，将所有空间结构压缩为单一向量，丢失了精密任务所需的细粒度空间信息——策略知道它\emph{看到了什么}，却不知道交互\emph{发生在哪里}。

我们提出\textbf{\fasafull{}（\fasa{}）}，一个通过帧间特征时间对比自监督发现任务相关空间区域的即插即用模块。\fasa{}的核心洞察是：连续帧之间变化最显著的区域对应于当前活跃的交互区域。基于这一洞察，\fasa{}将空间特征分解为动态流（交互区域）和静态流（稳定上下文），编码变化质心和运动方向，并提取空间关键点坐标——提供比全局池化丰富$2\times$的空间信息，无需任何外部教师模型。结合特征差分编码的"语义速度"，\fasa{}为流匹配策略提供了丰富的空间感知能力。

\fasa{}不需要人工空间标注，无需外部教师模型，并能无缝集成到现有流匹配架构中。我们将其集成到VITA流匹配策略中，在一个具有挑战性的6自由度机器人S型焊缝跟踪任务上进行评估，实验表明\fasa{}显著提升了跟踪精度。
\end{abstract}

% === 引言 ===
\section{引言}
\label{sec:introduction}

从专家示范中进行模仿学习已成为获取复杂机器人操作技能的主流范式~\cite{chi2023diffusion,zhao2023act,brohan2023rt2}。在生成式方法中，\emph{流匹配}~\cite{lipman2023flow,tong2024improving}作为扩散策略的替代方案近年来受到广泛关注。通过学习一个确定性速度场，将简单先验分布通过常微分方程（ODE）传输到目标动作分布，流匹配具有多项优势：无需噪声调度的稳定训练、快速确定性采样，以及结合最优传输（OT-CFM）~\cite{tong2024improving}时的直线传输路径。VITA框架~\cite{gao2026vita}是这一方法的典型代表，证明了从视觉观测隐变量直接流向动作隐变量可以获得优于条件扩散策略的多模态覆盖能力。

尽管取得了上述进展，当前的流匹配策略——以及大多数模仿学习架构——在视觉编码阶段存在一个重要的局限性：

\paragraph{空间信息丢失。}标准观测编码器对卷积特征图施加全局平均池化，将所有空间结构压缩为单个向量。这丢弃了关于任务相关物体\emph{在视野中位置}的关键信息。对于焊缝跟踪等精密任务，策略不仅需要知道它\emph{看到了什么}，还需要空间布局——焊缝位置、接近方向以及工具与工件之间的相对几何关系。全局池化将一个$512 \times H \times W$的丰富特征图压缩为512维向量，空间分辨率从数百个位置降为零。

在本文中，我们提出\textbf{\fasafull{}（\fasa{}）}，一个通过帧间特征时间对比自监督发现任务相关空间区域的即插即用空间注意力模块（图~\ref{fig:architecture}）：

\begin{itemize}[leftmargin=*,itemsep=4pt]
    \item \textbf{自监督空间注意力}（第~\ref{sec:spatial}节）：\fasa{}的核心洞察是\textbf{连续帧之间变化最显著的区域对应于当前活跃的交互区域}。通过检测帧间特征变化，\fasa{}将空间特征分解为动态（交互）流和静态（上下文）流，编码变化质心和方向，并提取空间关键点坐标——提供比全局池化丰富$2\times$的空间信息，无需任何外部教师模型。

    \item \textbf{语义速度编码}（第~\ref{sec:temporal_diff}节）：与简单堆叠连续帧不同，我们计算显式的特征空间差分$\Delta \bm{f}_t = \bm{f}_t - \bm{f}_{t-1}$，编码特征空间中的方向性变化。这种"语义速度"在接近阶段大（场景快速变化）、跟踪阶段小（细微偏移），提供了隐式的阶段信号。
\end{itemize}

\fasa{}是\emph{完全自监督的}：它不需要空间标注或训练时的外部教师模型。\fasa{}的变化精炼网络通过策略的动作预测损失端到端训练，自动学习抑制无关变化（如光照）并放大任务相关变化（如工具-工件接触）。

我们将\fasa{}集成到VITA流匹配架构中，在一个具有挑战性的6自由度S型焊缝跟踪任务上进行评估，其中安川MR2010机械臂需要在仿真环境中接近随机放置的工件并精确跟随弯曲的焊接轨迹。实验表明，相比基础VITA策略和多个基线方法，\fasa{}显著提升了跟踪精度。

\paragraph{贡献。}本文的主要贡献如下：
\begin{enumerate}[leftmargin=*,itemsep=2pt]
    \item 我们识别出空间信息丢失是流匹配视觉运动策略的重要局限性，并提出\fasa{}（Flow-Aware Spatial Attention）作为即插即用的自监督解决方案。
    \item \fasa{}通过多统计量帧间特征时间对比自监督发现任务相关区域，将空间特征分解为动态流和静态流，编码变化质心和方向，并提取空间关键点坐标，消除了对外部教师模型的依赖。
    \item 结合特征差分的"语义速度"编码，\fasa{}仅需连续观测帧即可在标准卷积特征图上操作，使其广泛适用于本文研究的特定架构和任务之外。
    \item 我们在精密机器人跟踪任务上验证了\fasa{}的有效性，全面的消融研究证实了空间感知的关键贡献。
\end{enumerate}

% === 相关工作 ===
\section{相关工作}
\label{sec:related_work}

\subsection{面向机器人策略的流匹配}

流匹配~\cite{lipman2023flow,liu2023flow}学习一个时间相关的速度场$v_\theta(x_t, t)$，通过ODE $\mathrm{d}x_t = v_\theta(x_t, t)\,\mathrm{d}t$将样本从先验分布传输到数据分布。与扩散模型~\cite{ho2020denoising,song2021score}相比，流匹配具有更简单的训练过程（无需噪声调度）、确定性采样，以及与最优传输（OT-CFM）~\cite{tong2024improving}结合时更直的传输路径。

在机器人领域，Diffusion Policy~\cite{chi2023diffusion}首次证明去噪扩散可以有效建模视觉运动控制的多模态动作分布，引入了动作分块（action chunking）范式。VITA~\cite{gao2026vita}将这一思想扩展到流匹配，提出从视觉观测隐变量直接流向动作隐变量——消除了随机噪声起始点，实现了自然的多模态动作生成。$\pi_0$~\cite{black2024pi0}将流匹配策略扩展到基础模型级别，通过视觉-语言-动作预训练展示了跨任务的强零样本泛化能力。然而，所有这些方法在视觉编码阶段都采用全局平均池化，丢失了精密任务所需的空间结构信息。本文的工作建立在VITA架构基础之上，通过自监督空间注意力对其进行增强。

\subsection{观测堆叠与帧差分}

最简单的时间编码策略是\emph{观测堆叠}：将最近$k$帧拼接作为输入~\cite{mnih2015human}。虽然直接简单，但这仅提供局部时间上下文，未显式编码帧间\emph{变化内容}，且随时间窗长度增加扩展性差。我们的特征差分方法（第~\ref{sec:temporal_diff}节）通过计算显式特征空间差分来解决这一问题，在不随观测窗增加维度的情况下编码"语义速度"。这种差分表示也为\fasa{}的变化检测提供了自然的基础。

\subsection{视觉运动策略中的空间注意力}

空间注意力机制旨在引导视觉编码器关注任务相关区域，而非平等对待所有空间位置。

\paragraph{学习的空间表示。}Spatial Softmax~\cite{levine2016end}从卷积特征图计算逐通道的期望坐标，提供紧凑的空间表示，已成为视觉运动策略学习的标准工具。然而，它对所有通道施加均匀注意力，不区分任务相关和无关区域。

\paragraph{通道与空间注意力模块。}Squeeze-and-Excitation（SE）网络~\cite{hu2018squeeze}通过全局池化和轻量级MLP学习通道注意力权重，按重要性重新加权特征通道。CBAM~\cite{woo2018cbam}通过额外的空间注意力分支学习二维注意力图进行扩展。Coordinate Attention~\cite{hou2021coordinate}将通道注意力分解为两个空间方向，捕获沿各轴的长程位置依赖关系。虽然这些模块对一般视觉任务有效，但它们从数据分布中学习\emph{静态}注意力模式，不会适应episode内当前的交互动态。

\paragraph{基于基础模型的知识蒸馏。}视觉基础模型如Grounding DINO~\cite{liu2023grounding}和CLIP~\cite{radford2021learning}编码了丰富的空间和语义知识，可通过注意力热图匹配或特征蒸馏损失传递给轻量级主干网络。虽然有效，但这些方法引入了对特定教师模型的依赖，需要离线特征提取流水线，且产生不随任务阶段变化的静态空间目标。

我们的\fasa{}模块采取了根本不同的方法：它不依赖外部空间监督或学习静态注意力模式，而是通过帧间特征时间对比\emph{自监督}地发现任务相关区域。核心洞察是，连续帧之间视觉变化最大的区域对应于活跃的交互区——恰恰是操作任务中最相关的区域。这使得\fasa{}的注意力天然具有\emph{动态}性，随着任务从接近阶段到跟踪阶段实时适应。

% === 预备知识 ===
\section{预备知识}
\label{sec:preliminaries}

\subsection{最优传输条件流匹配}

流匹配~\cite{lipman2023flow}学习一个时间相关的速度场$v_\theta: \RR^d \times [0,1] \to \RR^d$，生成一条概率路径$p_t$，将先验分布$p_0$（如标准高斯分布）插值到数据分布$p_1$。样本$x_0 \sim p_0$通过积分ODE传输到$x_1 \sim p_1$：
\begin{equation}
    \frac{\mathrm{d}x_t}{\mathrm{d}t} = v_\theta(x_t, t), \quad x_0 \sim p_0, \quad t \in [0, 1].
    \label{eq:flow_ode}
\end{equation}

条件流匹配（CFM）目标通过回归条件速度场$u_t(x | x_1)$来训练$v_\theta$，该条件速度场生成条件概率路径$p_t(x | x_1)$：
\begin{equation}
    \LL_{\text{CFM}} = \EE_{t \sim \mathcal{U}[0,1],\, x_1 \sim p_1,\, x_t \sim p_t(\cdot|x_1)} \left\| v_\theta(x_t, t) - u_t(x_t | x_1) \right\|^2.
    \label{eq:cfm_loss}
\end{equation}

使用最优传输CFM（OT-CFM）~\cite{tong2024improving}时，条件路径为具有OT耦合的线性插值：
\begin{equation}
    x_t = (1 - t)\, x_0 + t\, x_1, \quad u_t(x_t | x_1) = x_1 - x_0,
    \label{eq:ot_interpolation}
\end{equation}
其中$(x_0, x_1)$通过OT计划$\pi^* = \argmin_\pi \EE_{(x_0,x_1)\sim\pi}\|x_0 - x_1\|^2$进行匹配。这产生比标准CFM更直的流轨迹和更快的收敛。

在推理时，通过$K$步欧拉积分求解ODE生成动作：
\begin{equation}
    x_{t+\Delta t} = x_t + v_\theta(x_t, t) \cdot \Delta t, \quad \Delta t = 1/K.
    \label{eq:euler_sampling}
\end{equation}

\subsection{VITA架构}
\label{sec:vita_arch}

VITA~\cite{gao2026vita}以隐空间公式实例化面向视觉运动策略的流匹配。架构由四个组件组成（图~\ref{fig:architecture}，灰色区域）：

\paragraph{观测编码器$\FF_\phi$。}ResNet-18~\cite{he2016deep}主干网络从$N$个相机图像中提取视觉特征，经过全局平均池化后与本体感受状态$s_t \in \RR^{d_s}$拼接：
\begin{equation}
    \bm{o}_t = \FF_\phi\!\left(\{I_t^{(i)}\}_{i=1}^N,\; s_t\right) \in \RR^{d_o}, \quad d_o = 512N + d_s.
    \label{eq:obs_encoder}
\end{equation}
线性投影将其映射到隐空间：$\bm{z}_t^{\text{obs}} = W_{\text{enc}}\, \bm{o}_t \in \RR^{d_z}$。

\paragraph{动作自编码器。}一维CNN编码器$E_\psi$将动作块$\bm{a}_{t:t+H} \in \RR^{H \times d_a}$映射到隐向量$\bm{z}^{\text{act}} = E_\psi(\bm{a}_{t:t+H}) \in \RR^{d_z}$，MLP解码器$D_\psi$重建动作：$\hat{\bm{a}}_{t:t+H} = D_\psi(\bm{z}^{\text{act}})$。

\paragraph{流网络$G_\theta$。}带有FiLM~\cite{perez2018film}时间条件化的MLP预测速度场。每层$l$执行：
\begin{equation}
    \bm{h}^{(l)} = \bm{h}^{(l-1)} + \gamma_t^{(l)} \odot \text{MLP}\!\left((\bm{1} + \bm{\alpha}_t^{(l)}) \odot \text{LN}(\bm{h}^{(l-1)}) + \bm{\beta}_t^{(l)}\right),
    \label{eq:film_layer}
\end{equation}
其中$(\gamma_t, \bm{\alpha}_t, \bm{\beta}_t) = f^{(l)}(\text{SinEmb}(t))$是从流时间步$t$的正弦嵌入导出的FiLM参数。

\paragraph{训练目标。}VITA使用复合损失：
\begin{equation}
    \LL_{\text{VITA}} = \LL_{\text{flow}} + \lambda_{\text{enc}} \LL_{\text{enc}} + \lambda_{\text{fld}} \LL_{\text{fld}} + \lambda_{\text{con}} \LL_{\text{con}},
    \label{eq:vita_loss}
\end{equation}
其中$\LL_{\text{flow}}$是OT-CFM损失（公式~\ref{eq:cfm_loss}），$\LL_{\text{enc}}$是动作自编码器重建损失，$\LL_{\text{fld}}$是流解码重建损失，$\LL_{\text{con}} = \|\bm{z}^{\text{act}}_{\text{pred}} - \bm{z}^{\text{act}}\|^2$是隐空间一致性损失。关键的是，流从观测隐变量$\bm{z}_t^{\text{obs}}$而非随机噪声开始，建立了直接的视觉到动作映射。

\paragraph{空间信息瓶颈。}在上述公式中，每个相机图像的特征图$\bm{F}^{(i)} \in \RR^{512 \times H' \times W'}$经过全局平均池化压缩为单个512维向量。这意味着所有空间结构——物体位置、交互区域、几何关系——在信息进入隐空间之前就已丢失。本文提出的\fasa{}正是替代这一池化操作，在保持即插即用兼容性的同时恢复丰富的空间信息。

% === 方法 ===
\section{方法}
\label{sec:method}

\subsection{概述}
\label{sec:method_overview}

标准流匹配视觉运动策略采用全局平均池化将卷积特征图压缩为单一向量，丢失了精密任务所需的空间结构信息。我们提出\textbf{\fasafull{}（\fasa{}）}，一个即插即用的自监督空间注意力模块，通过两个互补机制解决这一问题：

\begin{enumerate}[leftmargin=*,itemsep=4pt]
    \item \textbf{基于特征差分的语义速度}（第~\ref{sec:temporal_diff}节）：通过计算连续帧特征的显式差分，编码特征空间中的方向性变化，为空间注意力提供时间对比基础。

    \item \textbf{\fasa{}空间注意力}（第~\ref{sec:spatial}节）：通过检测帧间特征变化自监督发现交互区域，将空间特征分解为动态流和静态流，编码变化质心和方向，提取空间关键点坐标——提供比全局池化丰富$2\times$的空间信息，无需外部教师模型或空间标注。
\end{enumerate}

\fasa{}完全自监督且即插即用：其变化精炼网络通过策略的动作预测损失端到端训练，无需专门的辅助损失。它仅需连续观测帧并在标准卷积特征图上操作，可与任何CNN主干网络配合使用。图~\ref{fig:architecture}提供了架构概览。

% ==================================================================
% 语义速度
% ==================================================================
\subsection{基于特征差分的语义速度}
\label{sec:temporal_diff}

与其简单地堆叠连续帧$[\bm{f}_{t-1}, \bm{f}_t]$（这会使维度加倍但未编码\emph{变化内容}），我们计算显式的特征空间差分：
\begin{equation}
    \bm{o}_t = [\bm{f}_t;\; \underbrace{\bm{f}_t - \bm{f}_{t-1}}_{\text{"语义速度"}};\; \bm{s}_t;\; \bm{s}_t - \bm{s}_{t-1}] \in \RR^{2(d_f N + d_s)},
    \label{eq:temporal_obs}
\end{equation}
其中$\bm{f}_t$是逐相机视觉特征（第~\ref{sec:spatial}节），$\bm{s}_t$是本体感受状态，$N$是相机数量。特征差分$\Delta \bm{f}_t = \bm{f}_t - \bm{f}_{t-1}$编码了特征空间中的方向性变化——接近时大（场景快速变化），跟踪时小（细微偏移）——提供了隐式的阶段信号。这种"语义速度"表示具有固定维度$2(d_f N + d_s)$，不随观测窗长度变化。

语义速度与\fasa{}互补：\fasa{}提供丰富的\emph{逐帧}空间特征，而特征差分编码\emph{跨帧}的变化方向和幅度。两者共同使策略既知道交互"发生在哪里"，也知道交互"如何变化"。

% ==================================================================
% 核心：空间感知
% ==================================================================
\subsection{自监督空间感知：\fasa{}}
\label{sec:spatial}

全局平均池化将$C \times H \times W$的特征图压缩为单个$C$维向量，丢失所有空间信息。对于精密任务，策略需要知道交互\emph{发生在哪里}。我们提出\fasa{}（\fasafull{}），通过一个简单但强大的洞察恢复空间信息：\textbf{连续帧之间变化最大的区域是与当前任务最相关的区域}。

% --- 变化检测 ---
\subsubsection{多统计量变化检测}
\label{sec:fasa_change}

给定来自ResNet主干网络的连续特征图$\bm{F}_t, \bm{F}_{t-1} \in \RR^{C \times H \times W}$，\fasa{}计算逐元素绝对差分$\bm{\Delta} = |\bm{F}_t - \bm{F}_{t-1}|$并提取三个互补统计量：
\begin{equation}
    \bm{C}_{\text{in}} = \left[\;
        \underbrace{\textstyle\text{mean}_c\, \bm{\Delta}}_{\text{平均强度}}\;;\;
        \underbrace{\max_c \bm{\Delta}}_{\text{峰值响应}}\;;\;
        \underbrace{\text{std}_c\, \bm{\Delta}}_{\text{通道一致性}}
    \;\right] \in \RR^{3 \times H \times W}.
    \label{eq:change_input}
\end{equation}

一个可学习的精炼网络$R_\rho$（4层CNN：$3 \to 32 \to 32 \to 1$，带批归一化和ReLU）将这些统计量处理为归一化的变化注意力图，其互补构成稳定注意力图：
\begin{equation}
    \bm{A}_{\text{change}} = \text{normalize}\!\left(\sigma(R_\rho(\bm{C}_{\text{in}}))\right), \quad
    \bm{A}_{\text{stable}} = \text{normalize}\!\left(1 - \sigma(R_\rho(\bm{C}_{\text{in}}))\right),
    \label{eq:change_attn}
\end{equation}
其中归一化确保空间权重之和为一。精炼网络$R_\rho$学习抑制无关变化（如光照）并放大任务相关变化（如工具-工件接触），通过策略的动作预测损失端到端训练。

% --- 空间特征提取 ---
\subsubsection{丰富的空间特征提取}
\label{sec:fasa_features}

从注意力图出发，\fasa{}为每个相机提取四种互补的空间表示：

\textbf{（1）动态特征}（$\RR^{512}$）：交互区域上的注意力加权池化：
\begin{equation}
    \bm{f}_{\text{dyn}} = \textstyle\sum_{y,x} \bm{A}_{\text{change}}(y,x) \cdot \bm{F}_t(:,y,x).
    \label{eq:dynamic_pool}
\end{equation}

\textbf{（2）静态特征}（$\RR^{512}$）：稳定上下文上的注意力加权池化：
\begin{equation}
    \bm{f}_{\text{sta}} = \textstyle\sum_{y,x} \bm{A}_{\text{stable}}(y,x) \cdot \bm{F}_t(:,y,x).
    \label{eq:static_pool}
\end{equation}

\textbf{（3）变化质心和方向}（$\RR^4$）：交互发生的位置及其移动方向：
\begin{equation}
    [c_x, c_y, d_x, d_y], \quad (c_x, c_y) = \textstyle\sum_{y,x} \bm{A}_{\text{change}} \cdot (x,y), \quad (d_x, d_y) = \text{centroid}_t - \text{centroid}_{t-1}.
    \label{eq:change_coords}
\end{equation}

\textbf{（4）空间Softmax坐标}（$\RR^{d_{\text{coord}}}$）：逐通道的期望位置~\cite{levine2016end}，投影至$d_{\text{coord}} = 32$维。这些提供了独立于变化检测的细粒度关键点信息。

统一的逐相机输出为：
\begin{equation}
    \bm{f}_{\fasa{}} = [\bm{f}_{\text{dyn}};\; \bm{f}_{\text{sta}};\; c_x, c_y, d_x, d_y;\; \bm{f}_{\text{coord}}] \in \RR^{1060},
    \label{eq:fasa_output}
\end{equation}
替代512维的全局平均池化，提供$2.07\times$更丰富的空间信息。

\textbf{优雅降级。}当没有前一帧可用时（第一步），\fasa{}对两个流均回退到全局平均池化，变化坐标置零，确保输出维度一致。

\textbf{设计优势。}与基于蒸馏的空间注意力不同，\fasa{}是（1）\emph{自监督的}——不需要教师模型；（2）\emph{即插即用的}——给定连续帧即可与任何CNN主干网络配合使用；（3）\emph{动态的}——注意力实时适应当前交互；（4）\emph{无辅助损失的}——完全通过策略的主损失函数端到端训练。

% ==================================================================
% 训练与推理
% ==================================================================
\subsection{训练与推理}
\label{sec:training}

\subsubsection{训练目标}

\method{}框架端到端训练（包括ResNet主干网络），使用与基础VITA相同的损失函数：
\begin{equation}
    \LL_{\method{}} = \LL_{\text{flow}} + \lambda_{\text{enc}}\,\LL_{\text{enc}} + \lambda_{\text{fld}}\,\LL_{\text{fld}} + \lambda_{\text{con}}\,\LL_{\text{con}}.
    \label{eq:total_loss}
\end{equation}
\fasa{}没有专门的辅助损失——其变化精炼网络完全通过动作预测损失的梯度学习。这是\fasa{}的一个重要设计特点：它不引入可能干扰主要优化目标的额外损失项，而是自然地从流匹配损失中学习任务相关的空间注意力。

\subsubsection{可选的主干网络预蒸馏}

主干网络可以通过离线蒸馏从Grounding DINO~\cite{liu2023grounding}可选地初始化，传递空间注意力（对检测热图的BCE损失）和语义特征（对Swin-T特征的余弦相似度损失）。蒸馏头随后被丢弃，主干网络继续端到端微调。这提供了更好的起始点但不是必需的——\fasa{}可以从零开始学习有效的空间特征。

\subsubsection{推理}

算法~\ref{alg:inference}总结了推理流程。在每个重规划步骤：（1）\fasa{}从连续帧提取空间特征；（2）特征差分计算语义速度；（3）观测编码器将丰富的空间特征映射到隐空间；（4）隐变量通过ODE积分流向动作隐变量；（5）动作解码器产生一个动作块，其中$H_a$个动作在重规划前被执行。

\begin{algorithm}[t]
\caption{\method{} 推理}
\label{alg:inference}
\begin{algorithmic}[1]
\REQUIRE 训练好的策略$\pi_{\method{}}$，episode长度$L$
\STATE 初始化队列$\mathcal{Q} \leftarrow \emptyset$，$\bm{F}_{-1} \leftarrow \texttt{None}$
\FOR{$t = 0, 1, \ldots, L-1$}
    \IF{$\mathcal{Q}$为空}
        \STATE $\bm{F}_t \leftarrow \text{ResNet}(\{I_t^{(i)}\})$ \hfill \textit{// 主干网络特征图}
        \STATE $\bm{f}_t \leftarrow \fasa{}(\bm{F}_t, \bm{F}_{t-1})$ \hfill \textit{// 空间感知（每相机1060D）}
        \STATE $\bm{o}_t \leftarrow [\bm{f}_t;\, \Delta\bm{f}_t;\, s_t;\, \Delta s_t]$ \hfill \textit{// 语义速度}
        \STATE $\bm{z}_t^{\text{obs}} \leftarrow W_{\text{enc}}\, \bm{o}_t$ \hfill \textit{// 观测编码}
        \STATE $\bm{a}_{t:t+H} \leftarrow D_\psi\!\left(\text{ODE}(G_\theta,\, \bm{z}_t^{\text{obs}})\right)$ \hfill \textit{// 流 + 解码}
        \STATE 入队$\bm{a}_{t:t+H_a}$；\; $\bm{F}_{t-1} \leftarrow \bm{F}_t$
    \ENDIF
    \STATE 执行$a_t \leftarrow \mathcal{Q}.\text{dequeue}()$
\ENDFOR
\end{algorithmic}
\end{algorithm}

% === 实验 ===
\section{实验}
\label{sec:experiments}

\subsection{任务描述}
\label{sec:task}

我们在仿真环境中的6自由度机器人S型焊缝跟踪任务上评估\method{}（图~\ref{fig:task}）。安川MR2010工业机械臂需要执行一个两阶段任务：

\begin{enumerate}[leftmargin=*,itemsep=2pt]
    \item \textbf{接近阶段}（约200步）：机器人从中性位姿移动到随机放置工件上弯曲焊缝的起始点。该阶段需要快速、定向的运动和粗略定位。
    \item \textbf{跟踪阶段}（约500步）：机器人精确跟随工件表面的S型焊缝轨迹。该阶段要求缓慢、精确的运动，达到亚毫米级精度。
\end{enumerate}

工件位置在不同episode间随机化，要求策略在不同几何配置间泛化。任务在MuJoCo~\cite{todorov2012mujoco}中仿真，具有真实的关节动力学和接触建模。

\subsection{数据集}
\label{sec:dataset}

我们使用带有解析逆运动学的脚本控制器收集了160条专家示范轨迹。每条轨迹包含约700个时间步，以50~Hz录制。观测空间包括：
\begin{itemize}[leftmargin=*,itemsep=1pt]
    \item 两个RGB相机图像（$480 \times 640$）：俯视图和末端执行器安装相机；
    \item 6维关节角度（本体感受状态）。
\end{itemize}
动作空间为6维关节位置指令（绝对值）。数据集共计112,000帧，以LeRobot~\cite{cadene2024lerobot}格式存储，并附带逐特征归一化统计量。

\subsection{基线方法}
\label{sec:baselines}

我们将\method{}与以下方法进行比较：

\begin{itemize}[leftmargin=*,itemsep=2pt]
    \item \textbf{VITA}~\cite{gao2026vita}：基础流匹配策略，单帧观测（$\text{obs\_horizon}=1$），全局平均池化。
    \item \textbf{VITA + obs\_horizon=2}：使用朴素双帧堆叠的VITA（特征拼接但不做差分）。
    \item \textbf{VITA + 特征差分}：使用$\text{obs\_horizon}=2$和特征差分（语义速度）的VITA，但仍使用全局平均池化，无\fasa{}。
    \item \textbf{ACT}~\cite{zhao2023act}：Action Chunking with Transformers，基于Transformer的模仿学习基线。
\end{itemize}

\subsection{评估指标}
\label{sec:metrics}

我们对每种方法在50个随机工件位置的episode上进行评估，报告以下指标：

\begin{itemize}[leftmargin=*,itemsep=1pt]
    \item \textbf{接近成功率}（\%）：末端执行器到达焊缝起始点5mm以内的episode比例。
    \item \textbf{跟踪距离}（mm）：跟踪阶段末端执行器与真实焊缝最近点之间的平均、中位数和最大欧氏距离。
    \item \textbf{多级覆盖率}（\%）：跟踪步骤中距离低于0.5mm、1mm、2mm和5mm阈值的比例。
    \item \textbf{奖励等级分布}：被评为S（$<$0.5mm平均）/ A（$<$1mm）/ B（$<$2mm）/ C（$<$5mm）/ D（$<$10mm）/ E（$<$20mm）/ FAIL（$\geq$20mm）的episode比例。
\end{itemize}

\subsection{实现细节}
\label{sec:impl_details}

所有模型使用隐空间维度$d_z = 512$，4层流网络（MLP比率4.0），动作块大小$H = 32$，动作窗口$H_a = 4$。流采样使用$K = 6$步欧拉积分。训练100k梯度步，使用AdamW优化器（$\beta_1 = 0.95$，$\beta_2 = 0.999$，权重衰减$10^{-6}$），批大小32，余弦学习率调度含500步预热。推理时使用指数$0.75$的EMA。ResNet-18主干网络与策略端到端训练（不冻结）。

基础VITA损失权重为$\lambda_{\text{enc}} = 0.5$，$\lambda_{\text{fld}} = 0.5$，$\lambda_{\text{con}} = 1.0$。

对于\fasa{}：每个相机的\fasa{}模块产生1060维特征向量（$512 + 512 + 4 + 32$），替代标准的512维全局平均池化输出。变化精炼网络$R_\rho$是4层CNN（$3 \to 32 \to 32 \to 1$），带批归一化，输入3通道（通道级特征差分的均值、最大值、标准差）。坐标投影每个相机增加约33k参数。\fasa{}要求$\text{obs\_horizon} \geq 2$以提供变化检测所需的连续帧；在每个episode的第一步，优雅降级为全局池化加空间softmax。

对于视觉预蒸馏（可选）：使用Grounding DINO的Swin-T主干网络（来自~\cite{liu2023grounding}的预训练权重），为两个相机的所有112,000帧预提取空间注意力热图和768维语义特征。预蒸馏阶段对ResNet主干网络训练40k步，使用注意力BCE损失和特征余弦相似度损失。预蒸馏后，蒸馏头被丢弃，主干网络在策略训练期间继续端到端微调。

基础VITA模型有约31.9M参数；加入\fasa{}后增加到约32.4M（每个相机仅增加约67k参数）。所有实验在单块NVIDIA RTX 4080 SUPER GPU（32GB）上运行。

\subsection{主要结果}
\label{sec:main_results}

表~\ref{tab:main_results}给出了主要比较结果。\textit{结果将在实验完成后填入。}

\begin{table*}[t]
\centering
\caption{S型焊缝跟踪任务的主要结果。最优结果用\textbf{粗体}标注。所有指标为50个评估episode的平均值。}
\label{tab:main_results}
\begin{tabular}{lccccccc}
\toprule
\multirow{2}{*}{\textbf{方法}} & \textbf{接近} & \multicolumn{3}{c}{\textbf{跟踪距离 (mm)}} & \multicolumn{3}{c}{\textbf{覆盖率 (\%)}} \\
\cmidrule(lr){3-5} \cmidrule(lr){6-8}
 & \textbf{成功率 (\%)} & 平均 & 中位数 & 最大 & $<$1mm & $<$2mm & $<$5mm \\
\midrule
ACT~\cite{zhao2023act}            & -- & -- & -- & -- & -- & -- & -- \\
VITA~\cite{gao2026vita}           & -- & -- & -- & -- & -- & -- & -- \\
VITA + obs\_horizon=2             & -- & -- & -- & -- & -- & -- & -- \\
VITA + 特征差分            & -- & -- & -- & -- & -- & -- & -- \\
\midrule
\method{}（本文）                  & -- & -- & -- & -- & -- & -- & -- \\
\bottomrule
\end{tabular}
\end{table*}

\subsection{消融实验}
\label{sec:ablation}

为理解\fasa{}各组件的贡献，我们进行了系统的消融实验（表~\ref{tab:ablation}）。我们评估以下配置：

\begin{enumerate}[label=(\alph*),leftmargin=*,itemsep=1pt]
    \item VITA基线（全局平均池化，$\text{obs\_horizon}=1$）
    \item VITA + obs\_horizon=2（朴素双帧堆叠）
    \item VITA + 特征差分（语义速度，无\fasa{}）
    \item \fasa{}仅动态流（无静态流、无质心编码）
    \item \fasa{}无变化精炼网络（原始差分作为注意力）
    \item 完整\fasa{}（动态+静态+质心+坐标）
\end{enumerate}

\begin{table*}[t]
\centering
\caption{消融实验。TD = 特征差分（语义速度），Dyn = 动态流，Sta = 静态流，CC = 变化质心/方向，SSM = 空间Softmax坐标，Ref = 变化精炼网络。}
\label{tab:ablation}
\begin{tabular}{cccccc|ccccc}
\toprule
\textbf{TD} & \textbf{Dyn} & \textbf{Sta} & \textbf{CC} & \textbf{SSM} & \textbf{Ref} & \textbf{接近 (\%)} & \textbf{平均 (mm)} & \textbf{中位数 (mm)} & \textbf{$<$2mm (\%)} & \textbf{$<$5mm (\%)} \\
\midrule
 &  &  &  &  &  & -- & -- & -- & -- & -- \\
\checkmark &  &  &  &  &  & -- & -- & -- & -- & -- \\
\checkmark & \checkmark &  &  &  &  & -- & -- & -- & -- & -- \\
\checkmark & \checkmark &  &  &  & \checkmark & -- & -- & -- & -- & -- \\
\checkmark & \checkmark & \checkmark & \checkmark & \checkmark &  & -- & -- & -- & -- & -- \\
\checkmark & \checkmark & \checkmark & \checkmark & \checkmark & \checkmark & -- & -- & -- & -- & -- \\
\bottomrule
\end{tabular}
\end{table*}

\subsection{分析}
\label{sec:analysis}

\paragraph{\fasa{}注意力可视化。}图~\ref{fig:fasa_attention}可视化了\fasa{}在不同任务阶段产生的变化注意力图$\bm{A}_{\text{change}}$。在接近阶段，我们期望看到分散的注意力，因为整个场景快速变化。在跟踪阶段，注意力应集中在工具尖端附近的焊缝区域，证明\fasa{}在没有任何显式空间监督的情况下学会了关注任务相关区域。\textit{实验完成后将添加分析。}

\paragraph{动态流 vs. 静态流。}\fasa{}将空间特征分解为动态流（交互区域特征）和静态流（稳定上下文特征）。我们分析两者在不同任务阶段的行为差异：接近阶段动态流变化剧烈，而跟踪阶段两者趋于一致。\textit{实验完成后将添加分析。}

\paragraph{泛化能力。}我们通过在比训练集更大分布范围内采样的工件位置上进行测试来评估泛化能力。\textit{实验完成后将添加结果。}

% === 讨论 ===
\section{讨论}
\label{sec:discussion}

\paragraph{空间感知为何重要。}全局平均池化是视觉运动策略中最常见的特征聚合方式，但它将$512 \times H' \times W'$的丰富特征图压缩为单个512维向量。对于粗粒度任务（如抓取大型物体），这种信息丢失可能无关紧要。但对于精密任务（如焊缝跟踪、插入、缝合），空间信息至关重要——策略需要知道工具尖端在哪里、焊缝在哪里、两者之间的相对位置关系如何。\fasa{}通过自监督变化检测恢复了这些空间信息，且仅增加极少的参数开销。

\paragraph{自监督 vs. 基于蒸馏的空间注意力。}视觉基础模型蒸馏（如从Grounding DINO或CLIP）可以提供空间注意力监督，但存在重要局限：（1）需要访问特定的教师模型和离线特征提取流水线；（2）蒸馏的注意力是\emph{静态的}——反映教师的通用目标检测能力，而非特定任务的动态交互模式；（3）不能跨领域泛化，除非重新运行教师模型。\fasa{}通过从任务本身的时间结构中推导空间注意力来解决所有三个局限。基于变化的注意力自然适应当前相关的内容：接近时，它突出显示在静态场景中移动的机器人；跟踪时，它聚焦于微妙的工具-工件接触区域。在我们的框架中，可选的预蒸馏增强了主干网络的特征质量（有利于\fasa{}的变化检测），而\fasa{}在池化层面提供动态的、任务自适应的空间注意力。

\paragraph{焊缝跟踪之外的适用性。}虽然我们在焊缝跟踪任务上评估\fasa{}，但其核心机制——帧间特征变化检测并据此分配空间注意力——是完全任务无关的。任何需要细粒度空间理解的精密任务——如插入、缝合或接触密集型操作——都可以从\fasa{}中受益。\fasa{}的自监督特性意味着迁移到新领域时不需要任务特定的标注或重新训练教师模型。

\paragraph{计算开销。}表~\ref{tab:compute}总结了参数和推理时间开销。\fasa{}模块每个相机仅增加约67k参数（两个相机共约134k），就提供了$2.07\times$更丰富的空间特征（1060D vs. 512D），是一个极其高效的升级。相比基础VITA的约31.9M参数，\fasa{}仅增加约0.4\%的参数量。

\begin{table}[t]
\centering
\caption{\fasa{}相对于基础VITA策略的计算开销。}
\label{tab:compute}
\begin{tabular}{lcc}
\toprule
\textbf{组件} & \textbf{参数量} & \textbf{推理时间 (ms)} \\
\midrule
基础VITA（全局平均池化） & $\sim$31.9M & -- \\
\midrule
\multicolumn{3}{l}{\textit{\fasa{}模块}} \\
\quad 变化精炼网络（每相机） & +$\sim$34k & +$<$0.3 \\
\quad 坐标投影（每相机） & +$\sim$33k & +$<$0.1 \\
\quad obs\_encoder 维度增大 & +$\sim$0.5M & +$<$0.1 \\
\midrule
\textbf{完整\method{}} & $\sim$32.4M (+1.6\%) & +$<$0.5 \\
\bottomrule
\end{tabular}
\end{table}

\paragraph{局限性。}以下局限性需要说明：
\begin{itemize}[leftmargin=*,itemsep=2pt]
    \item \textbf{\fasa{}冷启动。}\fasa{}需要$\text{obs\_horizon} \geq 2$来计算帧间变化。在每个episode的第一步（无前一帧），它降级为全局池化加空间softmax。对于第一个动作至关重要的任务，这可能影响初始性能。
    \item \textbf{依赖视觉变化。}\fasa{}的注意力基于帧间特征变化。在极慢运动或静止阶段，变化信号可能过弱，导致注意力退化为近似均匀分布。不过在实践中，50~Hz采样下即使缓慢运动也产生可检测的特征变化。
    \item \textbf{单任务评估。}我们在单个任务领域上进行评估。虽然\fasa{}设计为通用的即插即用模块，但在多样化操作任务上的更广泛评估将加强我们的主张。
\end{itemize}

\paragraph{未来方向。}有前景的扩展包括：（1）在真实机器人硬件上部署\method{}以验证空间感知的仿真到真实迁移；（2）探索在多个ResNet层输出上操作的多尺度\fasa{}以实现层次化空间注意力；（3）将\fasa{}的变化注意力图与好奇心驱动探索结合用于强化学习；（4）将\fasa{}扩展到多帧时间对比（超越两个连续帧）以实现更丰富的运动模式检测；（5）探索将时间感知通过非视觉通道（如关节速度、动作历史）安全地注入策略，避免干扰视觉特征学习。

% === 结论 ===
\section{结论}
\label{sec:conclusion}

我们提出了\fasa{}（\fasafull{}），一个通过帧间特征时间对比自监督发现任务相关空间区域的即插即用空间注意力模块。\fasa{}的核心洞察简单而强大：连续帧之间变化最显著的区域对应于当前活跃的交互区域。基于这一洞察，\fasa{}将空间特征分解为动态流（交互区域）和静态流（稳定上下文），编码变化质心和运动方向，并提取空间关键点坐标——提供比全局平均池化丰富$2.07\times$的空间信息，每个相机仅增加约67k参数。

\fasa{}完全自监督，无需空间标注、外部教师模型或专门的辅助损失。它仅需连续观测帧并在标准卷积特征图上操作，使其广泛适用于本文研究的特定架构和任务之外。结合特征差分编码的"语义速度"，\fasa{}为流匹配策略提供了丰富的空间感知和隐式的时间变化感知。

我们将\fasa{}集成到VITA流匹配架构中，在具有挑战性的6自由度S型焊缝跟踪任务上进行评估。实验表明\fasa{}显著提升了跟踪精度，消融实验证实了各空间特征组件（动态池化、静态池化、变化质心编码、空间softmax坐标）的互补贡献。

我们认为自监督空间注意力是机器人模仿学习中广泛重要的能力。\fasa{}所采取的方法——从观测流本身提取空间注意力，而非依赖人工标注或外部教师模型——代表了构建更强大、更便携、更具泛化能力的机器人策略的有前景的方向。


\bibliographystyle{unsrtnat}
\bibliography{../paper/references}

\end{document}
